\documentclass[12pt]{article}

\usepackage{fontspec}
\setmainfont{Times New Roman}
\usepackage[a4paper,margin=2.5cm]{geometry}
\usepackage{setspace}
\onehalfspacing
\usepackage{longtable,booktabs,array,calc}
\usepackage{caption}
\newcounter{none}
\usepackage{fancyvrb}
\usepackage{graphicx}
\usepackage{hyperref}
\hypersetup{
  colorlinks=true,
  linkcolor=black,
  urlcolor=blue,
  citecolor=black
}
\IfFileExists{xurl.sty}{\usepackage{xurl}}{}
\setlength{\parindent}{1.25cm}
\setlength{\parskip}{0pt}
\providecommand{\tightlist}{%
  \setlength{\itemsep}{0pt}\setlength{\parskip}{0pt}}
\newlength{\cslhangindent}
\setlength{\cslhangindent}{1.5em}
\newlength{\csllabelwidth}
\setlength{\csllabelwidth}{3em}
\newenvironment{CSLReferences}[2]
 {\begin{list}{}{%
  \setlength{\itemindent}{0pt}
  \setlength{\leftmargin}{0pt}
  \setlength{\parsep}{0pt}
  \ifodd #1
    \setlength{\leftmargin}{\cslhangindent}
    \setlength{\itemindent}{-\cslhangindent}
  \fi
  \setlength{\itemsep}{#2\baselineskip}}}
 {\end{list}}
\newcommand{\CSLBlock}[1]{\hfill\break#1\hfill\break}
\newcommand{\CSLLeftMargin}[1]{\parbox[t]{\csllabelwidth}{#1}}
\newcommand{\CSLRightInline}[1]{\parbox[t]{\linewidth - \csllabelwidth}{#1}\break}
\newcommand{\CSLIndent}[1]{\hspace{\cslhangindent}#1}
\newcommand{\citeproctext}{}

\title{Portal Nacional de Contratações Públicas e governo aberto: uma
análise exploratória das capitais do Sudeste\\\large Transparência
formal, dados abertos e fragmentação institucional}
\author{PREENCHER\_NOME}
\date{2026-07-01}

\begin{document}
\maketitle

\begin{center}
Universidade de São Paulo\\
Governo Aberto\\
Docente: Jorge Machado
\end{center}

\thispagestyle{empty}
\newpage
\setcounter{page}{1}

\section{Introdução}\label{introduuxe7uxe3o}

Este trabalho analisa o Portal Nacional de Contratações Públicas (PNCP)
como infraestrutura de governo aberto nas capitais do Sudeste
brasileiro. A pergunta de pesquisa é: em que medida a publicação de
contratações no PNCP favorece transparência, reutilização de dados e
controle social nas prefeituras de São Paulo, Rio de Janeiro, Belo
Horizonte e Vitória?

O tema dialoga com a disciplina Governo Aberto porque combina acesso à
informação, padrões tecnológicos, dados abertos e accountability. A
própria disciplina orienta que o trabalho articule problema, objetivos,
método e conclusões, com atenção à estrutura textual e às normas de
citação (Wikiversidade, 2026). O PNCP é especialmente relevante porque a
Lei nº 14.133/2021 o institui como ambiente nacional de divulgação de
contratações públicas (Brasil, 2021).

O argumento central é que o PNCP amplia a transparência formal ao
centralizar metadados, documentos e links de origem, mas a efetividade
como governo aberto depende da completude dos registros e da capacidade
de reconstituir a organização administrativa de cada prefeitura. O caso
de São Paulo mostra essa tensão com mais nitidez: grande parte dos
registros municipais aparece fora do CNPJ matriz do município.

\section{Objetivos e método}\label{objetivos-e-muxe9todo}

O objetivo geral é avaliar, de forma exploratória, como o PNCP expressa
princípios de governo aberto nas contratações das capitais do Sudeste.
Os objetivos específicos são: comparar volume e concentração
institucional dos registros, verificar a completude de campos
essenciais, observar a presença de documentos e discutir a fragmentação
de CNPJs em São Paulo.

O recorte empírico considera pregões eletrônicos, modalidade 6 no PNCP,
publicados entre 2026-01-01 e 2026-05-31. Foram usados os CNPJs matriz
das prefeituras de Rio de Janeiro, Belo Horizonte e Vitória. Para São
Paulo, combinou-se o CNPJ matriz com uma varredura por UF e código IBGE,
filtrando entidades municipais executivas.

{\def\LTcaptype{none} % do not increment counter
\begin{longtable}[]{@{}llll@{}}
\toprule\noalign{}
Capital & Fonte de coleta & Registros brutos & Limite \\
\midrule\noalign{}
\endhead
\bottomrule\noalign{}
\endlastfoot
Sao Paulo & matrix\_cnpj & 6 & \\
Sao Paulo & municipality\_scan & 250 & 5 \\
Rio de Janeiro & matrix\_cnpj & 347 & \\
Belo Horizonte & matrix\_cnpj & 90 & \\
Vitoria & matrix\_cnpj & 97 & \\
\end{longtable}
}

A amostra final selecionou até dez contratações por capital, por sorteio
pseudoaleatório reprodutível após ordenação por
\texttt{numeroControlePNCP}. A semente usada foi 20260608. A opção por
amostra fixa evita que capitais com mais registros dominem a
interpretação.

\section{Referencial teórico}\label{referencial-teuxf3rico}

Governo aberto costuma ser descrito como um arranjo que combina
transparência, participação social, colaboração e prestação de contas. A
Open Government Partnership associa o tema a compromissos de abertura,
integridade pública e participação cidadã (Open Government Partnership,
2011). A OCDE também define governo aberto como uma cultura de
governança que promove princípios de transparência, integridade,
accountability e participação das partes interessadas (Organisation for
Economic Co-operation and Development, 2017).

Dados abertos são uma dimensão operacional desse debate. Para que a
publicação seja útil, não basta disponibilizar informação em páginas
isoladas: é necessário permitir acesso, reuso, padronização e combinação
com outras bases (W3C Brasil, 2011). A Lei de Acesso à Informação também
reforça a publicidade como regra e o sigilo como exceção (Brasil, 2011).
Assim, transparência pública deve ser avaliada não apenas pela
existência de dados, mas por sua usabilidade para fiscalização e
pesquisa.

No caso brasileiro, o PNCP constitui uma infraestrutura relevante para
esse tipo de análise, pois centraliza informações e documentos de
contratações públicas (Brasil, 2026c). A existência de API e de dados
consultáveis amplia o potencial de reutilização por cidadãos,
pesquisadores e jornalistas (Brasil, 2026b, 2026a).

\section{Desenvolvimento e análise}\label{desenvolvimento-e-anuxe1lise}

\subsection{Visão comparativa}\label{visuxe3o-comparativa}

{\def\LTcaptype{none} % do not increment counter
\begin{longtable}[]{@{}
  >{\raggedright\arraybackslash}p{(\linewidth - 10\tabcolsep) * \real{0.1667}}
  >{\raggedright\arraybackslash}p{(\linewidth - 10\tabcolsep) * \real{0.1667}}
  >{\raggedright\arraybackslash}p{(\linewidth - 10\tabcolsep) * \real{0.1667}}
  >{\raggedright\arraybackslash}p{(\linewidth - 10\tabcolsep) * \real{0.1667}}
  >{\raggedright\arraybackslash}p{(\linewidth - 10\tabcolsep) * \real{0.1667}}
  >{\raggedright\arraybackslash}p{(\linewidth - 10\tabcolsep) * \real{0.1667}}@{}}
\toprule\noalign{}
\begin{minipage}[b]{\linewidth}\raggedright
Capital
\end{minipage} & \begin{minipage}[b]{\linewidth}\raggedright
Candidatos
\end{minipage} & \begin{minipage}[b]{\linewidth}\raggedright
Amostra
\end{minipage} & \begin{minipage}[b]{\linewidth}\raggedright
CNPJs distintos
\end{minipage} & \begin{minipage}[b]{\linewidth}\raggedright
Registros na matriz
\end{minipage} & \begin{minipage}[b]{\linewidth}\raggedright
Participação da matriz
\end{minipage} \\
\midrule\noalign{}
\endhead
\bottomrule\noalign{}
\endlastfoot
Sao Paulo & 61 & 10 & 10 & 6 & 9.8\% \\
Rio de Janeiro & 347 & 10 & 1 & 347 & 100.0\% \\
Belo Horizonte & 90 & 10 & 1 & 90 & 100.0\% \\
Vitoria & 97 & 10 & 1 & 97 & 100.0\% \\
\end{longtable}
}

A primeira diferença relevante está na concentração institucional. Rio
de Janeiro, Belo Horizonte e Vitória aparecem concentradas no CNPJ
matriz no recorte adotado. São Paulo, por outro lado, apresenta
múltiplos CNPJs municipais executivos associados a unidades como fundo
municipal, hospital, companhia de engenharia de tráfego, secretaria e
subprefeitura.

\subsection{Fragmentação em São
Paulo}\label{fragmentauxe7uxe3o-em-suxe3o-paulo}

{\def\LTcaptype{none} % do not increment counter
\begin{longtable}[]{@{}ll@{}}
\toprule\noalign{}
Métrica & Valor \\
\midrule\noalign{}
\endhead
\bottomrule\noalign{}
\endlastfoot
CNPJ matriz & 46395000000139 \\
Registros candidatos & 61 \\
Registros no CNPJ matriz & 6 \\
Registros fora do CNPJ matriz & 55 \\
Participação fora da matriz & 90.2\% \\
CNPJs distintos & 10 \\
CNPJs fora da matriz & 9 \\
\end{longtable}
}

Esse resultado é substantivo para governo aberto. Uma consulta limitada
ao CNPJ matriz de São Paulo recuperaria apenas uma fração dos registros
elegíveis. Portanto, a transparência formal do PNCP convive com uma
barreira de reconstrução institucional: o cidadão precisa saber que a
prefeitura pode publicar contratações por diferentes CNPJs municipais.

\subsection{Completude e
rastreabilidade}\label{completude-e-rastreabilidade}

{\def\LTcaptype{none} % do not increment counter
\begin{longtable}[]{@{}lllll@{}}
\toprule\noalign{}
Capital & Campo & Presentes & Amostra & Percentual \\
\midrule\noalign{}
\endhead
\bottomrule\noalign{}
\endlastfoot
Sao Paulo & Valor homologado & 7 & 10 & 70.0\% \\
Sao Paulo & Link de origem & 10 & 10 & 100.0\% \\
Sao Paulo & Documentos & 10 & 10 & 100.0\% \\
Rio de Janeiro & Valor homologado & 6 & 10 & 60.0\% \\
Rio de Janeiro & Link de origem & 10 & 10 & 100.0\% \\
Rio de Janeiro & Documentos & 10 & 10 & 100.0\% \\
Belo Horizonte & Valor homologado & 7 & 10 & 70.0\% \\
Belo Horizonte & Link de origem & 9 & 10 & 90.0\% \\
Belo Horizonte & Documentos & 10 & 10 & 100.0\% \\
Vitoria & Valor homologado & 3 & 10 & 30.0\% \\
Vitoria & Link de origem & 0 & 10 & 0.0\% \\
Vitoria & Documentos & 10 & 10 & 100.0\% \\
\end{longtable}
}

Em todos os municípios da amostra, objeto, datas básicas, unidade
administrativa e documentos estiveram amplamente disponíveis. A
principal variação apareceu em valor homologado e link do sistema de
origem. Vitória, por exemplo, teve documentos em todos os itens
amostrados, mas nenhum dos dez registros trouxe
\texttt{linkSistemaOrigem}, o que reduz a rastreabilidade para o
ambiente de origem.

{\def\LTcaptype{none} % do not increment counter
\begin{longtable}[]{@{}llllll@{}}
\toprule\noalign{}
Capital & Com documentos & Amostra & Mín. & Máx. & Média \\
\midrule\noalign{}
\endhead
\bottomrule\noalign{}
\endlastfoot
Sao Paulo & 10 & 10 & 1 & 1 & 1.0 \\
Rio de Janeiro & 10 & 10 & 1 & 1 & 1.0 \\
Belo Horizonte & 10 & 10 & 1 & 1 & 1.0 \\
Vitoria & 10 & 10 & 2 & 12 & 3.8 \\
\end{longtable}
}

\subsection{Constatações
empíricas}\label{constatauxe7uxf5es-empuxedricas}

\begin{itemize}
\tightlist
\item
  Na amostra candidata de Sao Paulo, 90.2\% dos registros elegiveis
  ficaram fora do CNPJ matriz; nas demais capitais analisadas, os
  registros por CNPJ matriz concentraram 100\% dos candidatos coletados.
\item
  Sao Paulo apresentou 10 CNPJs distintos no recorte elegivel, enquanto
  Rio de Janeiro, Belo Horizonte e Vitoria apareceram com um CNPJ cada
  no recorte por matriz.
\item
  Vitoria teve documentos vinculados em todos os itens da amostra, mas
  nenhum dos 10 registros amostrados trouxe \texttt{linkSistemaOrigem},
  o que reduz a rastreabilidade para o sistema de origem.
\item
  O valor homologado apareceu com menor completude em Vitoria (30.0\%),
  indicando que parte dos registros estava em fase anterior ao resultado
  ou sem homologacao registrada no recorte.
\item
  A quantidade de documentos anexados variou de forma relevante: em
  Vitoria, um item da amostra chegou a 12 documentos, enquanto outros
  municipios tiveram padrao mais concentrado.
\end{itemize}

\subsection{Exemplos de registros}\label{exemplos-de-registros}

Os exemplos completos foram deslocados para o Apêndice A. No corpo do
texto, basta notar que um registro da Secretaria Municipal de Saúde de
São Paulo retorna CNPJ do Fundo Municipal de Saúde, unidade
\texttt{PMSP\ -\ SECRETARIA\ MUNICIPAL\ DE\ SAÚDE}, objeto de aquisição
hospitalar e documento de edital. Já o exemplo no CNPJ matriz retorna
\texttt{MUNICIPIO\ DE\ SAO\ PAULO} e unidade da Secretaria do Governo
Municipal. Essa diferença concreta sustenta a leitura de fragmentação
institucional.

\section{Conclusões}\label{conclusuxf5es}

A análise indica que o PNCP é uma infraestrutura importante de governo
aberto: ele centraliza registros, expõe campos padronizados, oferece
documentos e permite automação por API. Esses elementos fortalecem a
transparência formal e criam condições para controle social.

Ao mesmo tempo, a comparação mostra que abertura não é sinônimo
automático de inteligibilidade. A fragmentação de CNPJs em São Paulo
torna a busca mais complexa e pode subestimar a atividade contratual
municipal quando o pesquisador consulta apenas o CNPJ matriz. A
contribuição do trabalho está em mostrar que a avaliação de governo
aberto deve observar também a arquitetura institucional dos dados.

Como conclusão regional exploratória, as capitais do Sudeste analisadas
apresentam boa disponibilidade básica de dados e documentos no PNCP, mas
diferem na forma de organização dos registros. A agenda de melhoria
passa por documentação mais clara dos CNPJs/unidades, maior completude
de links de origem e mecanismos que facilitem ao cidadão reconstruir o
universo de órgãos vinculados a cada prefeitura.

\section{Referências}\label{referuxeancias}

\protect\phantomsection\label{refs}
\begin{CSLReferences}{1}{0}
\bibitem[\citeproctext]{ref-lei12527}
Brasil. (2011). \emph{Lei nº 12.527, de 18 de novembro de 2011}.
\url{https://www.planalto.gov.br/ccivil_03/_ato2011-2014/2011/lei/l12527.htm}

\bibitem[\citeproctext]{ref-lei14133}
Brasil. (2021). \emph{Lei nº 14.133, de 1º de abril de 2021}.
\url{https://www.planalto.gov.br/ccivil_03/_ato2019-2022/2021/lei/l14133.htm}

\bibitem[\citeproctext]{ref-pncpApiConsulta}
Brasil. (2026a). \emph{API PNCP Consulta}.
\url{https://pncp.gov.br/pncp-consulta/v3/api-docs}

\bibitem[\citeproctext]{ref-pncpDadosAbertos}
Brasil. (2026b). \emph{PNCP: dados abertos}.
\url{https://www.gov.br/pncp/pt-br/acesso-a-informacao/dados-abertos}

\bibitem[\citeproctext]{ref-pncpPortal}
Brasil. (2026c). \emph{Portal Nacional de Contratações Públicas}.
\url{https://www.gov.br/pncp/pt-br/pncp}

\bibitem[\citeproctext]{ref-ogpDeclaration}
Open Government Partnership. (2011). \emph{Open Government Declaration}.
\url{https://www.opengovpartnership.org/process/joining-ogp/open-government-declaration/}

\bibitem[\citeproctext]{ref-oecd2017}
Organisation for Economic Co-operation and Development. (2017).
\emph{Recommendation of the Council on Open Government}.
\url{https://legalinstruments.oecd.org/en/instruments/OECD-LEGAL-0438}

\bibitem[\citeproctext]{ref-w3cManualDadosAbertos}
W3C Brasil. (2011). \emph{Manual dos dados abertos: governo}.
\url{https://www.w3c.br/pub/Materiais/PublicacoesW3C/Manual_Dados_Abertos_WEB.pdf}

\bibitem[\citeproctext]{ref-curso2026}
Wikiversidade. (2026). \emph{Governo Aberto (curso2026)}.
\url{https://pt.wikiversity.org/wiki/Governo_Aberto_(curso2026)}

\end{CSLReferences}

\appendix

\section{Apêndice A: exemplos compactos da
API}\label{apuxeandice-a-exemplos-compactos-da-api}

\subsection{Sao Paulo - exemplo fora do CNPJ
matriz}\label{sao-paulo---exemplo-fora-do-cnpj-matriz}

\begin{itemize}
\tightlist
\item
  Controle PNCP: \texttt{13864377000130-1-000007/2026}
\item
  CNPJ/órgão: \texttt{13864377000130} --- FUNDO MUNICIPAL DE SAUDE - FMS
\item
  Unidade: PMSP - SECRETARIA MUNICIPAL DE SAÚDE
\item
  Valor estimado: 863090.8
\item
  Valor homologado: 633080.0
\item
  Situação: Divulgada no PNCP
\item
  Documentos listados no exemplo: 1
\item
  Objeto: REGISTRO De PREÇOS OBJETIVANDO O FORNECIMENTO DE MATERIAIS DE
  OPME - SISTEMA DE PLACA E PARAFUSOPARA MÃO COM ENTREGA EM CONSIGNAÇÃO
  E COMODATO DE INSTRUMENTAIS E EQUIPAMENTOS, NECESSÁRIOS PARA O
  ATENDIMENTO DE CIRURGIAS NA ESPECIALIDADE DE ORTOPEDIA, A SEREM
  UTILIZADOS NAS UNIDADES HOSPITALARES PERTENCENTES À SECRETARIA
  MUNICIPAL DA SAÚDE DE SP, PARA O PERÍODO DE 12(DOZE) MESES.
\end{itemize}

\subsection{Sao Paulo - exemplo no CNPJ
matriz}\label{sao-paulo---exemplo-no-cnpj-matriz}

\begin{itemize}
\tightlist
\item
  Controle PNCP: \texttt{46395000000139-1-000028/2026}
\item
  CNPJ/órgão: \texttt{46395000000139} --- MUNICIPIO DE SAO PAULO
\item
  Unidade: PMSP - SECRETARIA DO GOVERNO MUNICIPAL
\item
  Valor estimado: 4652861.35
\item
  Valor homologado: 4652499.96
\item
  Situação: Divulgada no PNCP
\item
  Documentos listados no exemplo: 1
\item
  Objeto: Contratação de empresa especializada na prestação de serviços
  de controle de acesso do Edifício Conde Matarazzo, abrangendo a
  cobertura efetiva das portarias, a disponibilização de equipamentos de
  informática, a emissão de crachás e películas de identificação, bem
  como a manutenção preventiva e corretiva dos módulos de passagem,
  utilizando-se o software de propriedade da CONTRATANTE, com o
  fornecimento de peças necessárias. O objetivo é assegu\ldots{}
\end{itemize}

\subsection{Rio de Janeiro}\label{rio-de-janeiro}

\begin{itemize}
\tightlist
\item
  Controle PNCP: \texttt{42498733000148-1-000127/2026}
\item
  CNPJ/órgão: \texttt{42498733000148} --- MUNICIPIO DE RIO DE JANEIRO
\item
  Unidade: PREFEITURA MUNICIPAL DO RIO DE JANEIRO - RJ
\item
  Valor estimado: 99005.6
\item
  Valor homologado: 85452.38
\item
  Situação: Divulgada no PNCP
\item
  Documentos listados no exemplo: 1
\item
  Objeto: Registro de Preços para Aquisição de Materiais de Instalações
  Elétricas (C), para manutenção rotineira, a serem utilizados nas
  unidades físicas do Sistema BRT da MOBI-Rio.
\end{itemize}

\subsection{Belo Horizonte}\label{belo-horizonte}

\begin{itemize}
\tightlist
\item
  Controle PNCP: \texttt{18715383000140-1-000178/2026}
\item
  CNPJ/órgão: \texttt{18715383000140} --- MUNICIPIO DE BELO HORIZONTE
\item
  Unidade: SECRETARIA MUNICIPAL DE EDUCAÇÃO
\item
  Valor estimado: 111600.0
\item
  Valor homologado: 72000.0
\item
  Situação: Divulgada no PNCP
\item
  Documentos listados no exemplo: 1
\item
  Objeto: FORNECIMENTO DE AGUA MINERAL NATURAL EM GARRAFAO COM 20
  LITROS, INCLUINDO O EMPRESTIMO A TITULO DE COMODATO DE GARRAFAO COM
  CAPACIDADE PARA 20 LITROSVASILHAME, DESTINADOS AO ABASTECIMENTO DA
  SEDE DA SECRETARIA MUNICIPAL DE EDUCACAO ? SMED E GERENCIA REGIONAL DE
  EDUCACAO - BARREIRO , NOS TERMOS DA TABELA ABAIXO E CONFORME CONDICOES
  E EXIGENCIAS ESTABELECIDAS NESTE INSTRUMENTO.
\end{itemize}

\subsection{Vitoria}\label{vitoria}

\begin{itemize}
\tightlist
\item
  Controle PNCP: \texttt{27142058000126-1-000034/2026}
\item
  CNPJ/órgão: \texttt{27142058000126} --- MUNICIPIO DE VITORIA
\item
  Unidade: PREFEITURA MUNICIPAL DE VITORIA
\item
  Valor estimado: 669364.4
\item
  Valor homologado: não informado
\item
  Situação: Divulgada no PNCP
\item
  Documentos listados no exemplo: 2
\item
  Objeto: FORNECIMENTO DE ALIMENTOS (DESJEJUM, KIT LANCHE E MARMITEX).
\end{itemize}

\section{Apêndice B: declaração de uso de
IA}\label{apuxeandice-b-declarauxe7uxe3o-de-uso-de-ia}

Este trabalho utilizou ferramentas de inteligência artificial generativa
para apoio à programação do pipeline de coleta e análise, organização
dos resultados, revisão textual e geração inicial do relatório. A
seleção do tema, a validação das fontes, a interpretação dos resultados
e a versão final são de responsabilidade do(s) autor(es).

\section{Apêndice C: limitações
metodológicas}\label{apuxeandice-c-limitauxe7uxf5es-metodoluxf3gicas}

A pesquisa é exploratória e não representa todos os municípios do
Sudeste. A varredura municipal de São Paulo foi limitada
operacionalmente a cinco páginas da API para manter o pipeline
reprodutível em tempo razoável. Além disso, os resultados podem mudar
com novas publicações, retificações ou alterações na API do PNCP.

\end{document}
